\documentclass[10pt,a4paper]{amsart}
\usepackage{amsmath,amssymb,amsthm,graphicx}
% Retained minimal preamble for the Stacks Project chapter sites-cohomology.tex (Cohomology on Sites).
% Verbatim excerpt of the chapter preamble: Xy-pic setup (native preamble lines 14-18) and the native
% theorem declarations with the mathematical macros (native preamble lines 154-217). Copied mechanically;
% nothing added, removed or redefined. The wrapper class and the amsmath/amssymb/amsthm/graphicx base
% packages are supplied by the assembler, so the class/externaldocument part of the native preamble is omitted.
% For commutative diagrams we use Xy-pic
\usepackage[all]{xy}

% We use 2cell for 2-commutative diagrams.
\xyoption{2cell}
% Theorem environments.
%
\theoremstyle{plain}
\newtheorem{theorem}[subsection]{Theorem}
\newtheorem{proposition}[subsection]{Proposition}
\newtheorem{lemma}[subsection]{Lemma}

\theoremstyle{definition}
\newtheorem{definition}[subsection]{Definition}
\newtheorem{example}[subsection]{Example}
\newtheorem{exercise}[subsection]{Exercise}
\newtheorem{situation}[subsection]{Situation}

\theoremstyle{remark}
\newtheorem{remark}[subsection]{Remark}
\newtheorem{remarks}[subsection]{Remarks}

\numberwithin{equation}{subsection}

% Macros
%
\def\lim{\mathop{\mathrm{lim}}\nolimits}
\def\colim{\mathop{\mathrm{colim}}\nolimits}
\def\Spec{\mathop{\mathrm{Spec}}}
\def\Hom{\mathop{\mathrm{Hom}}\nolimits}
\def\Ext{\mathop{\mathrm{Ext}}\nolimits}
\def\SheafHom{\mathop{\mathcal{H}\!\mathit{om}}\nolimits}
\def\SheafExt{\mathop{\mathcal{E}\!\mathit{xt}}\nolimits}
\def\Sch{\mathit{Sch}}
\def\Mor{\mathop{\mathrm{Mor}}\nolimits}
\def\Ob{\mathop{\mathrm{Ob}}\nolimits}
\def\Sh{\mathop{\mathit{Sh}}\nolimits}
\def\NL{\mathop{N\!L}\nolimits}
\def\CH{\mathop{\mathrm{CH}}\nolimits}
\def\proetale{{pro\text{-}\acute{e}tale}}
\def\etale{{\acute{e}tale}}
\def\QCoh{\mathit{QCoh}}
\def\Ker{\mathop{\mathrm{Ker}}}
\def\Im{\mathop{\mathrm{Im}}}
\def\Coker{\mathop{\mathrm{Coker}}}
\def\Coim{\mathop{\mathrm{Coim}}}

% Boxtimes
%
\DeclareMathSymbol{\boxtimes}{\mathbin}{AMSa}{"02}

%
% Macros for moduli stacks/spaces
%
\def\QCohstack{\mathcal{QC}\!\mathit{oh}}
\def\Cohstack{\mathcal{C}\!\mathit{oh}}
\def\Spacesstack{\mathcal{S}\!\mathit{paces}}
\def\Quotfunctor{\mathrm{Quot}}
\def\Hilbfunctor{\mathrm{Hilb}}
\def\Curvesstack{\mathcal{C}\!\mathit{urves}}
\def\Polarizedstack{\mathcal{P}\!\mathit{olarized}}
\def\Complexesstack{\mathcal{C}\!\mathit{omplexes}}
% \Pic is the operator that assigns to X its picard group, usage \Pic(X)
% \Picardstack_{X/B} denotes the Picard stack of X over B
% \Picardfunctor_{X/B} denotes the Picard functor of X over B
\def\Pic{\mathop{\mathrm{Pic}}\nolimits}
\def\Picardstack{\mathcal{P}\!\mathit{ic}}
\def\Picardfunctor{\mathrm{Pic}}
\def\Deformationcategory{\mathcal{D}\!\mathit{ef}}
\usepackage{cleveref}
\crefname{originallemma}{original lemma}{original lemmas}
\Crefname{originallemma}{Original lemma}{Original lemmas}
\makeatletter
\newlabel{categories-lemma-ameliorate-morphism-fibred-categories}{{08NF}{source}{}{sourcepointer-1}{}}
\newlabel{categories-lemma-ameliorate-morphism-fibred-categories@cref}{{[originallemma][0][]08NF}{[1][0][]source}{}{}{}}
\makeatother
\crefname{originalreference}{original reference}{original references}
\Crefname{originalreference}{Original reference}{Original references}
\makeatletter
\newlabel{lemma-compute-left-derived-pi-shriek}{{08PJ}{source}{}{sourcepointer-2}{}}
\newlabel{lemma-compute-left-derived-pi-shriek@cref}{{[originalreference][0][]08PJ}{[1][0][]source}{}{}{}}
\makeatother
\crefname{originallemma}{original lemma}{original lemmas}
\Crefname{originallemma}{Original lemma}{Original lemmas}
\makeatletter
\newlabel{sites-lemma-have-left-adjoint}{{08NH}{source}{}{sourcepointer-3}{}}
\newlabel{sites-lemma-have-left-adjoint@cref}{{[originallemma][0][]08NH}{[1][0][]source}{}{}{}}
\makeatother
\crefname{originallemma}{original lemma}{original lemmas}
\Crefname{originallemma}{Original lemma}{Original lemmas}
\makeatletter
\newlabel{sites-lemma-when-shriek}{{00XR}{source}{}{sourcepointer-4}{}}
\newlabel{sites-lemma-when-shriek@cref}{{[originallemma][0][]00XR}{[1][0][]source}{}{}{}}
\makeatother
\crefname{originallemma}{original lemma}{original lemmas}
\Crefname{originallemma}{Original lemma}{Original lemmas}
\makeatletter
\newlabel{sites-modules-lemma-have-left-adjoint}{{08P3}{source}{}{sourcepointer-5}{}}
\newlabel{sites-modules-lemma-have-left-adjoint@cref}{{[originallemma][0][]08P3}{[1][0][]source}{}{}{}}
\makeatother
\crefname{originalreference}{original reference}{original references}
\Crefname{originalreference}{Original reference}{Original references}
\makeatletter
\newlabel{situation-morphism-fibred-categories}{{08PA}{source}{}{sourcepointer-6}{}}
\newlabel{situation-morphism-fibred-categories@cref}{{[originalreference][0][]08PA}{[1][0][]source}{}{}{}}
\makeatother
\crefname{originallemma}{original lemma}{original lemmas}
\Crefname{originallemma}{Original lemma}{Original lemmas}
\makeatletter
\newlabel{stacks-lemma-topology-inherited-functorial}{{06NW}{source}{}{sourcepointer-7}{}}
\newlabel{stacks-lemma-topology-inherited-functorial@cref}{{[originallemma][0][]06NW}{[1][0][]source}{}{}{}}
\makeatother
\title{The higher derived lower shriek of an abelian sheaf along a morphism of fibred categories is computed by the homology of the index category}
\author{The Stacks Project Authors}
\date{stacks-project:sites-cohomology.tex:local-snapshot-20261011; GNU FDL 1.2 or later}
\begin{document}
\maketitle
\noindent\textit{Editorial note.} All delivered mathematical text is the author's own, copied from the pinned Stacks Project chapter file sites-cohomology.tex; no statement or proof step was rewritten, reordered or removed.\par\medskip
\noindent\textit{Editorial note.} The counted proof is the author's complete argument: the proof is written out in full: the description of the index category and the presheaf, the amelioration of the morphism with the exactness of the first factor, the reduction to the fibred case, and the final identification of the index categories.\par\medskip
\noindent\textit{Editorial note.} The label lemma-compute-left-derived-g-shriek names the counted statement, and the proof environment immediately following it in the source is delivered verbatim as its proof.\par\medskip
\noindent\textit{Editorial note.} The counted statement is carried by the author's own lemma environment in the chapter source.\par\medskip
\noindent\textit{Editorial note.} The statement refers to the standing situation and to the index categories introduced in the proof; the seven cited inputs are kept as external pointers with their canonical Stacks tags.\par\medskip
\noindent\textit{Editorial note.} The preamble is the chapter's own retained minimal preamble (sources/stacks-sites-cohomology/retained-preamble.tex), a verbatim excerpt of the native preamble: the Xy-pic setup and the native theorem declarations with the mathematical macros.\par\medskip
\noindent\textit{Allowed standard prerequisites (editorial).}\begin{itemize}
\item Situation of a morphism of fibred categories and its lower shriek g\_!. Same chapter, the standing situation of the morphism of fibred categories and the adjoint pair (g\_!, g\textasciicircum{}*). Supplies the object whose derived functors are computed and the index category appearing in the formula.
\item Amelioration of a morphism of fibred categories. Categories, the lemma factoring u through an ameliorated morphism u'' with an exact u'. Supplies the factorisation u = u'' o u' used to split the computation into an exact part and the general case.
\item Exactness of g'\_! and the computation lemma for the fibred case. Modules on Sites, the exactness of the lower shriek for the ameliorated morphism, and Same chapter, the computation of L\_n pi\_! for a fibred category. Supplies Lg\_! = Lg''\_! after applying the exact g'\_!, to which the known computation lemma applies.
\end{itemize}\medskip
\ifdefined\hypertarget\hypertarget{sourcepointer-1}{}\fi
\noindent\textit{Source reference.} Original lemma 08NF: sources/stacks-chow/tags.txt, source lines 10988--10988. Prerequisite pointer invoked at proof line(s) 10709; it names a standard fact of a sibling chapter that the counted argument uses, and every step of the counted proof built on it is retained verbatim in the delivered body.\par\medskip
\ifdefined\hypertarget\hypertarget{sourcepointer-2}{}\fi
\noindent\textit{Source reference.} Original reference 08PJ: sources/stacks-sites-cohomology/pinned-original-source.tex, source lines 10615--10615. Prerequisite pointer invoked at proof line(s) 10721; it names a standard fact of the same chapter that the counted argument uses, and every step of the counted proof built on it is retained verbatim in the delivered body.\par\medskip
\ifdefined\hypertarget\hypertarget{sourcepointer-3}{}\fi
\noindent\textit{Source reference.} Original lemma 08NH: sources/stacks-chow/tags.txt, source lines 10990--10990. Prerequisite pointer invoked at proof line(s) 10718; it names a standard fact of a sibling chapter that the counted argument uses, and every step of the counted proof built on it is retained verbatim in the delivered body.\par\medskip
\ifdefined\hypertarget\hypertarget{sourcepointer-4}{}\fi
\noindent\textit{Source reference.} Original lemma 00XR: sources/stacks-chow/tags.txt, source lines 1229--1229. Prerequisite pointer invoked at proof line(s) 10731; it names a standard fact of a sibling chapter that the counted argument uses, and every step of the counted proof built on it is retained verbatim in the delivered body.\par\medskip
\ifdefined\hypertarget\hypertarget{sourcepointer-5}{}\fi
\noindent\textit{Source reference.} Original lemma 08P3: sources/stacks-chow/tags.txt, source lines 11011--11011. Prerequisite pointer invoked at proof line(s) 10712; it names a standard fact of a sibling chapter that the counted argument uses, and every step of the counted proof built on it is retained verbatim in the delivered body.\par\medskip
\ifdefined\hypertarget\hypertarget{sourcepointer-6}{}\fi
\noindent\textit{Source reference.} Original reference 08PA: sources/stacks-sites-cohomology/pinned-original-source.tex, source lines 9767--9767. Prerequisite pointer invoked at proof line(s) 10680; it names a standard fact of the same chapter that the counted argument uses, and every step of the counted proof built on it is retained verbatim in the delivered body.\par\medskip
\ifdefined\hypertarget\hypertarget{sourcepointer-7}{}\fi
\noindent\textit{Source reference.} Original lemma 06NW: sources/stacks-chow/tags.txt, source lines 8550--8550. Prerequisite pointer invoked at proof line(s) 10729; it names a standard fact of a sibling chapter that the counted argument uses, and every step of the counted proof built on it is retained verbatim in the delivered body.\par\medskip
\begin{lemma}
\label{lemma-compute-left-derived-g-shriek}
Assumptions and notation as in
Situation \ref{situation-morphism-fibred-categories}.
For an abelian sheaf $\mathcal{F}'$ on $\mathcal{C}'$ the sheaf
$L_ng_!(\mathcal{F}')$ is the sheaf associated to the presheaf
$$
U \longmapsto H_n(\mathcal{I}_U, \mathcal{F}'_U)
$$
For notation and restriction maps see proof.
\end{lemma}

\begin{proof}
Say $p(U) = V$. The category $\mathcal{I}_U$ is the category of pairs
$(U', \varphi)$ where $\varphi : U \to u(U')$ is a morphism of $\mathcal{C}$
with $p(\varphi) = \text{id}_V$, i.e., $\varphi$ is a morphism of the
fibre category $\mathcal{C}_V$. Morphisms
$(U'_1, \varphi_1) \to (U'_2, \varphi_2)$ are given by morphisms
$a : U'_1 \to U'_2$ of the fibre category $\mathcal{C}'_V$ such that
$\varphi_2 = u(a) \circ \varphi_1$. The presheaf $\mathcal{F}'_U$ sends
$(U', \varphi)$ to $\mathcal{F}'(U')$.
We will construct the restriction mappings below.

\medskip\noindent
Choose a factorization
$$
\xymatrix{
\mathcal{C}' \ar@<1ex>[r]^{u'} &
\mathcal{C}'' \ar[r]^{u''} \ar@<1ex>[l]^w & \mathcal{C}
}
$$
of $u$ as in
Categories, Lemma \ref{categories-lemma-ameliorate-morphism-fibred-categories}.
Then $g_! = g''_! \circ g'_!$ and similarly for derived functors.
On the other hand, the functor $g'_!$ is exact, see
Modules on Sites, Lemma \ref{sites-modules-lemma-have-left-adjoint}.
Thus we get $Lg_!(\mathcal{F}') = Lg''_!(\mathcal{F}'')$ where
$\mathcal{F}'' = g'_!\mathcal{F}'$. Note that
$\mathcal{F}'' = h^{-1}\mathcal{F}'$ where
$h : \Sh(\mathcal{C}'') \to \Sh(\mathcal{C}')$ is the morphism of topoi
associated to $w$, see
Sites, Lemma \ref{sites-lemma-have-left-adjoint}.
The functor $u''$ turns $\mathcal{C}''$ into a fibred category
over $\mathcal{C}$, hence
Lemma \ref{lemma-compute-left-derived-pi-shriek}
applies to the computation of $L_ng''_!$. The result follows as the
construction of $\mathcal{C}''$ in the proof of
Categories, Lemma \ref{categories-lemma-ameliorate-morphism-fibred-categories}
shows that the fibre category $\mathcal{C}''_U$ is equal to
$\mathcal{I}_U$. Moreover, $h^{-1}\mathcal{F}'|_{\mathcal{C}''_U}$
is given by the rule described above
(as $w$ is continuous and cocontinuous by
Stacks, Lemma \ref{stacks-lemma-topology-inherited-functorial}
so we may apply
Sites, Lemma \ref{sites-lemma-when-shriek}).
\end{proof}

\end{document}
